\documentclass[10pt,letterpaper]{article}
\usepackage{cornell-notes}

\title{Chapter 8: Random Permutation}
\author{}
\date{}

\setDocTitle{Chapter 8: Random Permutation}
\setDocSubtitle{Combinatorial Algorithms}
\setDocVersion{v1.0}
\setDocStatus{Study Companion}
\setDocOwner{Combinatorial Algorithms Cornell Notes}
\setDocAuthor{}
\setDocDate{\today}
\setDocSummary{Cornell-format study and review notes for Chapter 8: Random Permutation.}

\setCornellCollection{Combinatorial Algorithms Cornell Notes}
\setCornellUnitType{Chapter}
\setCornellUnitNumber{8}
\setCornellUnitTitle{Random Permutation}
\setCornellCompanionLabel{Study and review companion}

\hypersetup{pdftitle={Chapter 8: Random Permutation},pdfsubject={Cornell notes},pdfauthor={}}

\begin{document}
\maketitle
\begin{tcolorbox}[colback=Paper,colframe=Ink]{\small\bfseries\color{Teal}CORNELL NOTES}\par{\LARGE\bfseries Chapter 8: Random Permutation of $n$ Letters (RANPER)}\par\medskip\textbf{Topic:} In-place uniform shuffling\hfill\textbf{Date:} \today\par\textbf{Study purpose:} Prove uniformity of decreasing-range random exchanges.\end{tcolorbox}
\begin{tcolorbox}[colback=Teal!5,colframe=Teal,title=Essential question]Why does choosing one uniformly random remaining position for each output position make every permutation equally likely?\end{tcolorbox}
\begin{tcolorbox}[colback=white,colframe=Mist,title=Learning objectives]\begin{itemize}\item Execute the in-place exchange loop.\item Prove the $1/n!$ probability of each permutation.\item Distinguish setup and shuffle-only modes.\item Identify off-by-one and random-integer hazards.\item Interpret the cycle-count sample check.\end{itemize}\textbf{Prerequisites:} permutations, conditional probability, array swaps.\end{tcolorbox}
\section*{Cornell notes}
\begin{longtable}{@{}Q!{\color{Mist}\vrule width 1pt}A@{}}\cornellhead
\cue{What is the procedure?}{Optionally initialize $a_i=i$. For $m=1,\ldots,n$, choose $L$ uniformly from $\{m,m+1,\ldots,n\}$ and exchange $a_m$ with $a_L$. After iteration $m$, position $m$ is fixed and the unchosen items occupy the remaining suffix.}
\cue{Why is it uniform?}{For any prescribed permutation, the first required item is selected with probability $1/n$, the second from the remaining items with probability $1/(n-1)$, and so on. Thus
\[
\Pr\{\pi\}=\frac1n\frac1{n-1}\cdots\frac11=\frac1{n!}.
\]
Every permutation follows one unique sequence of required choices.}
\cue{What does SETUP control?}{When true, the routine first writes $1,2,\ldots,n$ and returns a random permutation of those labels. When false, it shuffles the caller's existing array. The latter applies the same uniform position permutation to arbitrary items.}
\cue{Invariant}{Before iteration $m$, positions $1,\ldots,m-1$ contain a uniformly chosen ordered sample without replacement, and the suffix contains exactly the remaining items. Choosing uniformly within the suffix preserves the invariant.}
\cue{Cost and storage}{There are $n$ bounded random choices and at most $n$ swaps. The shuffle is in place and uses constant auxiliary storage beyond the input/output array.}
\cue{What can go wrong?}{Choose from the inclusive suffix, not from the whole array at every step. Mapping a random real to an integer must never produce $n+1$ and must not bias endpoints. Duplicated input values are shuffled uniformly by position, though distinct output value sequences may then have different multiplicities.}
\cue{How is the sample checked?}{The sample compares the observed mean number of cycles in random permutations with the exact value $1+\frac12+\cdots+\frac1n$. Agreement is a statistical diagnostic, not a proof of uniformity; the product-of-probabilities argument supplies the proof.}
\cue{Retrieval practice}{\begin{enumerate}\item What is the choice range at step $m$?\item Why does the product equal $1/n!$?\item What does setup=false mean?\item What invariant holds for the suffix?\item Which integer-mapping bug causes bias?\end{enumerate}}
\end{longtable}
\begin{tcolorbox}[colback=Ink!4,colframe=Ink,title=Cornell summary]
RANPER builds a permutation from left to right. At position $m$, it uniformly chooses one of the items not yet fixed, swaps that item into place, and continues on the shorter suffix. A fixed final permutation requires one specific choice at each step, with probabilities $1/n,1/(n-1),\ldots,1$, so its probability is exactly $1/n!$. The procedure works either after initializing the identity labels or directly on caller-supplied data. It is in-place and linear in the number of positions. The central implementation risk is generating an unbiased integer in the correct inclusive suffix range; an incorrect range produces a shuffle that may look random while not being uniform.
\end{tcolorbox}
\end{document}
